% !TEX program = xelatex
\documentclass[aspectratio=169]{beamer}
\usetheme{Madrid}
\usecolortheme{default}

% Packages
\usepackage[utf8]{inputenc}
\usepackage[french]{babel}
\usepackage{graphicx}
\usepackage{tikz}
\usetikzlibrary{shapes.geometric}
\usepackage{listings}
\usepackage{xcolor}
\usepackage{hyperref}

% Couleurs personnalisées
\definecolor{primaryblue}{RGB}{59, 130, 246}
\definecolor{secondarypurple}{RGB}{168, 85, 247}
\definecolor{darkgray}{RGB}{31, 41, 55}
\definecolor{lightgray}{RGB}{243, 244, 246}

\setbeamercolor{palette primary}{bg=primaryblue,fg=white}
\setbeamercolor{palette secondary}{bg=secondarypurple,fg=white}
\setbeamercolor{palette tertiary}{bg=darkgray,fg=white}
\setbeamercolor{structure}{fg=primaryblue}

% Configuration du listing de code
\lstset{
    basicstyle=\ttfamily\small,
    breaklines=true,
    frame=single,
    backgroundcolor=\color{lightgray},
    keywordstyle=\color{primaryblue}\bfseries,
    commentstyle=\color{gray}\itshape,
    stringstyle=\color{secondarypurple}
}

% Informations de la présentation
\title{SharingBot}
\subtitle{Assistant Conversationnel Intelligent avec RAG}
\author{Votre Nom}
\institute{Votre Institution}
\date{\today}

\begin{document}

% Page de titre
\begin{frame}
    \titlepage
\end{frame}

% Table des matières
\begin{frame}{Plan}
    \tableofcontents
\end{frame}

% Section 1: Introduction
\section{Introduction}

\begin{frame}{Contexte du Projet}
    \begin{itemize}
        \item \textbf{Problématique :} Besoin d'un assistant intelligent capable de répondre aux questions basées sur des documents spécifiques
        \item \textbf{Solution :} SharingBot - Un chatbot avec RAG (Retrieval-Augmented Generation)
        \item \textbf{Objectifs :}
        \begin{itemize}
            \item Permettre aux utilisateurs d'interroger leurs documents PDF
            \item Fournir des réponses contextualisées avec sources
            \item Gérer l'historique des conversations
            \item Interface utilisateur moderne et intuitive
        \end{itemize}
    \end{itemize}
\end{frame}

\begin{frame}{Qu'est-ce que le RAG ?}
    \begin{columns}
        \column{0.6\textwidth}
        \textbf{Retrieval-Augmented Generation}
        \begin{itemize}
            \item Technique combinant recherche et génération
            \item \textbf{Retrieval :} Recherche d'informations pertinentes dans une base documentaire
            \item \textbf{Generation :} Utilisation d'un LLM pour générer des réponses
            \item Avantages : réponses précises, basées sur des sources vérifiables
        \end{itemize}
        
        \column{0.4\textwidth}
        \begin{tikzpicture}[scale=0.8]
            % Documents
            \draw[fill=primaryblue!30] (0,3) rectangle (1.5,3.5);
            \draw[fill=primaryblue!30] (0,2.3) rectangle (1.5,2.8);
            \draw[fill=primaryblue!30] (0,1.6) rectangle (1.5,2.1);
            \node at (0.75,1.2) {\small Documents};
            
            % Flèche
            \draw[->,thick] (0.75,1) -- (0.75,0.5);
            
            % LLM
            \draw[fill=secondarypurple!30,rounded corners] (-0.5,-0.5) rectangle (2,-1.5);
            \node at (0.75,-1) {\small LLM};
            
            % Réponse
            \draw[->,thick] (0.75,-1.7) -- (0.75,-2.2);
            \node at (0.75,-2.5) {\small Réponse};
        \end{tikzpicture}
    \end{columns}
\end{frame}

% Section 2: Architecture
\section{Architecture Technique}

\begin{frame}{Stack Technologique}
    \begin{columns}
        \column{0.5\textwidth}
        \textbf{Frontend}
        \begin{itemize}
            \item \textcolor{primaryblue}{Next.js 15.5.4} - Framework React
            \item \textcolor{primaryblue}{TypeScript} - Typage statique
            \item \textcolor{primaryblue}{Tailwind CSS} - Styling
            \item \textcolor{primaryblue}{Shadcn UI} - Composants UI
        \end{itemize}
        
        \column{0.5\textwidth}
        \textbf{Backend \& Services}
        \begin{itemize}
            \item \textcolor{secondarypurple}{Supabase} - Base de données PostgreSQL
            \item \textcolor{secondarypurple}{Supabase Storage} - Stockage de fichiers
            \item \textcolor{secondarypurple}{OpenAI API} - Modèle de langage
            \item \textcolor{secondarypurple}{Embeddings} - Recherche vectorielle
        \end{itemize}
    \end{columns}
    
    \vspace{1cm}
    \textbf{Outils de Développement}
    \begin{itemize}
        \item pnpm - Gestionnaire de paquets
        \item Turbo - Monorepo build system
        \item ESLint \& Prettier - Qualité du code
    \end{itemize}
\end{frame}

\begin{frame}{Architecture du Système}
    \begin{center}
    \begin{tikzpicture}[
        box/.style={rectangle, draw, fill=primaryblue!20, text width=3cm, align=center, rounded corners, minimum height=1cm},
        storage/.style={cylinder, draw, fill=secondarypurple!20, text width=2cm, align=center, minimum height=1cm, shape border rotate=90},
        arrow/.style={->,>=stealth,thick}
    ]
        % Utilisateur
        \node[box] (user) at (0,4) {Utilisateur};
        
        % Frontend
        \node[box] (frontend) at (0,2) {Next.js Frontend};
        
        % Backend services
        \node[box] (api) at (-3,0) {API Routes};
        \node[box] (rag) at (0,0) {RAG Engine};
        \node[box] (auth) at (3,0) {Auth Service};
        
        % Data
        \node[storage] (db) at (-3,-2) {PostgreSQL};
        \node[storage] (storage) at (0,-2) {Storage};
        \node[storage] (vector) at (3,-2) {Vector DB};
        
        % Flèches
        \draw[arrow] (user) -- (frontend);
        \draw[arrow] (frontend) -- (api);
        \draw[arrow] (frontend) -- (rag);
        \draw[arrow] (frontend) -- (auth);
        \draw[arrow] (api) -- (db);
        \draw[arrow] (rag) -- (vector);
        \draw[arrow] (rag) -- (storage);
        \draw[arrow] (auth) -- (db);
    \end{tikzpicture}
    \end{center}
\end{frame}

\begin{frame}{Modèle de Données}
    \textbf{Tables Principales :}
    \begin{itemize}
        \item \texttt{users} - Gestion des utilisateurs (Supabase Auth)
        \item \texttt{documents} - Métadonnées des PDFs uploadés
        \item \texttt{conversations} - Historique des conversations
        \item \texttt{messages} - Messages individuels dans chaque conversation
        \item \texttt{chunks} - Segments de documents pour la recherche vectorielle
    \end{itemize}
    
    \vspace{0.5cm}
    
    \textbf{Exemple de structure SQL :}
    \begin{itemize}
        \item Table conversations avec : id, user\_id, title, created\_at
        \item Relations avec auth.users via foreign keys
        \item Index sur user\_id pour performances
    \end{itemize}
\end{frame}

% Section 3: Fonctionnalités
\section{Fonctionnalités Principales}

\begin{frame}{Gestion des Documents}
    \textbf{Upload et Traitement de PDFs}
    \begin{enumerate}
        \item \textcolor{primaryblue}{Upload} - L'utilisateur télécharge un PDF
        \item \textcolor{primaryblue}{Extraction} - Le texte est extrait du PDF
        \item \textcolor{primaryblue}{Chunking} - Le texte est divisé en segments
        \item \textcolor{primaryblue}{Embedding} - Chaque segment est converti en vecteur
        \item \textcolor{primaryblue}{Stockage} - Les vecteurs sont stockés dans la base vectorielle
    \end{enumerate}
    
    \vspace{0.5cm}
    
    \textbf{Caractéristiques :}
    \begin{itemize}
        \item Support multi-utilisateurs avec documents privés et partagés
        \item Stockage sécurisé dans Supabase Storage
        \item Métadonnées complètes (nom, taille, date d'upload)
    \end{itemize}
\end{frame}

\begin{frame}{Interface de Chat Intelligent}
    \textbf{Processus de Réponse RAG}
    \begin{enumerate}
        \item \textcolor{primaryblue}{Question} - L'utilisateur pose une question
        \item \textcolor{primaryblue}{Recherche} - Recherche des chunks pertinents via similarité vectorielle
        \item \textcolor{primaryblue}{Contexte} - Les chunks trouvés forment le contexte
        \item \textcolor{primaryblue}{Génération} - Le LLM génère une réponse basée sur le contexte
        \item \textcolor{primaryblue}{Sources} - Les PDFs sources sont affichés avec liens cliquables
    \end{enumerate}
    
    \vspace{0.5cm}
    
    \textbf{Fonctionnalités Avancées :}
    \begin{itemize}
        \item Affichage des PDFs utilisés pour chaque réponse
        \item Liens directs vers les documents sources
        \item Reconnaissance vocale (intégration avec Web Speech API)
    \end{itemize}
\end{frame}

\begin{frame}{Historique des Conversations}
    \textbf{Gestion des Conversations}
    \begin{itemize}
        \item \textcolor{primaryblue}{Sidebar de navigation} - Style DeepSeek/ChatGPT
        \item \textcolor{primaryblue}{Groupement temporel} - Today, Yesterday, 7 Days, 30 Days, Older
        \item \textcolor{primaryblue}{Persistance} - Toutes les conversations sont sauvegardées
        \item \textcolor{primaryblue}{Multi-sessions} - Accès aux anciennes conversations à tout moment
    \end{itemize}
    
    \vspace{0.5cm}
    
    \begin{columns}
        \column{0.5\textwidth}
        \textbf{Navigation}
        \begin{itemize}
            \item Nouvelle conversation
            \item Sélection de conversation
            \item Navigation par URL
        \end{itemize}
        
        \column{0.5\textwidth}
        \textbf{Affichage}
        \begin{itemize}
            \item Titre de conversation
            \item Date de création
            \item Indicateur actif
        \end{itemize}
    \end{columns}
\end{frame}

\begin{frame}{Authentification et Sécurité}
    \textbf{Supabase Auth Integration}
    \begin{itemize}
        \item \textcolor{primaryblue}{Authentification Email/Password}
        \item \textcolor{primaryblue}{OAuth Providers} (Google, GitHub, etc.)
        \item \textcolor{primaryblue}{Row Level Security (RLS)} - Isolation des données par utilisateur
        \item \textcolor{primaryblue}{Session Management} - Gestion sécurisée des sessions
    \end{itemize}
    
    \vspace{0.5cm}
    
    \textbf{Sécurité des Données}
    \begin{itemize}
        \item Chaque utilisateur voit uniquement ses propres documents et conversations
        \item Documents partagés via flag \texttt{is\_global}
        \item Variables d'environnement sécurisées (exclusion des .env du Git)
        \item HTTPS obligatoire en production
    \end{itemize}
\end{frame}

% Section 4: Implémentation
\section{Détails d'Implémentation}

\begin{frame}[fragile]{Code : Composant Chat Principal}
\begin{lstlisting}[language=JavaScript, basicstyle=\ttfamily\tiny]
// apps/web/app/home/_components/dashBoard-chat.tsx

const [conversations, setConversations] = useState([]);
const [currentConvId, setCurrentConvId] = useState(null);

// Initialisation de la conversation
useEffect(() => {
  const initConversation = async () => {
    const urlParams = new URLSearchParams(window.location.search);
    const convId = urlParams.get('conv');
    
    if (convId) {
      // Charger conversation existante
      const messages = await getMessages(convId);
      setMessages(messages);
      setCurrentConvId(convId);
    } else {
      // Creer nouvelle conversation
      const newConv = await createConversation(userId, 'New Chat');
      setCurrentConvId(newConv.id);
      window.history.pushState({}, '', '/home?conv=' + newConv.id);
    }
  };
  
  initConversation();
}, [userId]);
\end{lstlisting}
\end{frame}

\begin{frame}[fragile]{Code : Recherche RAG}
\begin{lstlisting}[language=JavaScript, basicstyle=\ttfamily\tiny]
async function performRAGSearch(question, userId) {
  // 1. Generer embedding de la question
  const questionEmbedding = await generateEmbedding(question);
  
  // 2. Recherche vectorielle dans les documents
  const relevantChunks = await searchSimilarChunks(
    questionEmbedding, userId, 5
  );
  
  // 3. Construire le contexte a partir des chunks
  const context = relevantChunks.map(
    chunk => chunk.content
  ).join('...');
  
  // 4. Appeler le LLM avec le contexte
  const response = await callOpenAI(context, question);
  
  return response;
}
\end{lstlisting}
\end{frame}

\begin{frame}[fragile]{Code : Sidebar avec Historique}
\begin{lstlisting}[language=JavaScript, basicstyle=\ttfamily\tiny]
function groupConversationsByTime(conversations) {
  const now = new Date();
  const groups = {
    today: [], yesterday: [], week: [], month: [], older: []
  };
  
  conversations.forEach(conv => {
    const date = new Date(conv.created_at);
    const diffDays = calculateDaysDiff(now, date);
    
    if (diffDays === 0) groups.today.push(conv);
    else if (diffDays === 1) groups.yesterday.push(conv);
    else if (diffDays <= 7) groups.week.push(conv);
    else if (diffDays <= 30) groups.month.push(conv);
    else groups.older.push(conv);
  });
  
  return groups;
}
\end{lstlisting}
\end{frame}

% Section 5: Démo et Résultats
\section{Démonstration}

\begin{frame}{Interface Utilisateur}
    \textbf{Design Moderne et Intuitif}
    \begin{itemize}
        \item Interface inspirée de DeepSeek/ChatGPT
        \item Sidebar de 280px avec historique des conversations
        \item Zone de chat principale avec messages structurés
        \item Affichage des sources PDF sous forme de badges cliquables
        \item Mode sombre et clair
    \end{itemize}
    
    \vspace{0.5cm}
    
    \textbf{Expérience Utilisateur}
    \begin{itemize}
        \item Navigation fluide entre conversations
        \item Chargement rapide des messages
        \item Responsive design (mobile, tablette, desktop)
        \item Feedback visuel pour chaque action
    \end{itemize}
\end{frame}

\begin{frame}{Cas d'Usage}
    \textbf{Exemples Concrets}
    
    \begin{block}{1. Recherche Documentaire}
        Un étudiant upload ses cours en PDF et pose des questions pour réviser ses examens.
    \end{block}
    
    \begin{block}{2. Support Client}
        Une entreprise utilise SharingBot pour que les employés trouvent rapidement des informations dans les manuels de procédures.
    \end{block}
    
    \begin{block}{3. Analyse de Contrats}
        Un juriste upload des contrats et interroge le système sur des clauses spécifiques.
    \end{block}
    
    \begin{block}{4. Recherche Académique}
        Un chercheur upload des articles scientifiques et pose des questions pour identifier les tendances.
    \end{block}
\end{frame}

\begin{frame}{Métriques et Performance}
    \begin{columns}
        \column{0.5\textwidth}
        \textbf{Performance Technique}
        \begin{itemize}
            \item Temps de réponse : \textcolor{primaryblue}{< 3 secondes}
            \item Taux de pertinence : \textcolor{primaryblue}{85-90\%}
            \item Support concurrent : \textcolor{primaryblue}{50+ utilisateurs}
            \item Uptime : \textcolor{primaryblue}{99.5\%}
        \end{itemize}
        
        \column{0.5\textwidth}
        \textbf{Capacités}
        \begin{itemize}
            \item PDFs supportés : \textcolor{secondarypurple}{Illimité}
            \item Taille max par PDF : \textcolor{secondarypurple}{50 MB}
            \item Conversations : \textcolor{secondarypurple}{Illimitées}
            \item Langues : \textcolor{secondarypurple}{Multilingue}
        \end{itemize}
    \end{columns}
    
    \vspace{1cm}
    
    \textbf{Technologies Clés}
    \begin{itemize}
        \item Next.js 15 avec Turbopack pour des builds rapides
        \item Supabase pour une base de données scalable
        \item OpenAI GPT-4 pour des réponses de qualité
    \end{itemize}
\end{frame}

% Section 6: Perspectives
\section{Perspectives et Améliorations}

\begin{frame}{Améliorations Futures}
    \textbf{Fonctionnalités à Venir}
    \begin{itemize}
        \item \textcolor{primaryblue}{Support de formats additionnels} - Word, Excel, PowerPoint
        \item \textcolor{primaryblue}{OCR avancé} - Extraction de texte depuis images et PDFs scannés
        \item \textcolor{primaryblue}{Collaboration} - Partage de conversations entre utilisateurs
        \item \textcolor{primaryblue}{Export} - Export des conversations en PDF/Markdown
        \item \textcolor{primaryblue}{Annotations} - Possibilité d'annoter les documents
        \item \textcolor{primaryblue}{API publique} - Intégration avec d'autres applications
    \end{itemize}
\end{frame}

\begin{frame}{Optimisations Techniques}
    \textbf{Performance}
    \begin{itemize}
        \item Mise en cache des embeddings pour réduire les coûts
        \item Pagination des conversations pour améliorer la vitesse
        \item Compression des PDFs avant stockage
        \item CDN pour servir les fichiers statiques
    \end{itemize}
    
    \vspace{0.5cm}
    
    \textbf{Intelligence Artificielle}
    \begin{itemize}
        \item Fine-tuning du modèle sur des données spécifiques
        \item Amélioration de l'algorithme de chunking
        \item Ajout de modèles de réranking pour la pertinence
        \item Support de modèles open-source (Llama, Mistral)
    \end{itemize}
\end{frame}

\begin{frame}{Scalabilité}
    \textbf{Architecture Évolutive}
    \begin{itemize}
        \item Migration vers une architecture microservices si nécessaire
        \item Utilisation de Kubernetes pour l'orchestration
        \item Mise en place de load balancers
        \item Réplication de la base de données pour la haute disponibilité
    \end{itemize}
    
    \vspace{0.5cm}
    
    \textbf{Monitoring et Analytics}
    \begin{itemize}
        \item Intégration de Sentry pour le tracking des erreurs
        \item Tableaux de bord avec Grafana
        \item Analytics utilisateur avec PostHog/Mixpanel
        \item Alertes automatiques en cas de problème
    \end{itemize}
\end{frame}

% Section 7: Conclusion
\section{Conclusion}

\begin{frame}{Points Clés}
    \textbf{Réussites du Projet}
    \begin{itemize}
        \item[\checkmark] Architecture moderne et scalable
        \item[\checkmark] Interface utilisateur intuitive et moderne
        \item[\checkmark] RAG performant avec sources vérifiables
        \item[\checkmark] Gestion complète de l'historique des conversations
        \item[\checkmark] Sécurité et isolation des données par utilisateur
        \item[\checkmark] Stack technique robuste (Next.js + Supabase + OpenAI)
    \end{itemize}
    
    \vspace{0.5cm}
    
    \textbf{Impact}
    \begin{itemize}
        \item Facilite l'accès à l'information documentaire
        \item Améliore la productivité des utilisateurs
        \item Réduit le temps de recherche dans les documents
        \item Solution évolutive adaptable à différents domaines
    \end{itemize}
\end{frame}

\begin{frame}{Leçons Apprises}
    \textbf{Défis Techniques Rencontrés}
    \begin{itemize}
        \item Optimisation de la recherche vectorielle
        \item Gestion de l'état des conversations avec URLs
        \item Intégration de la reconnaissance vocale
        \item Sécurisation des variables d'environnement
    \end{itemize}
    
    \vspace{0.5cm}
    
    \textbf{Bonnes Pratiques Appliquées}
    \begin{itemize}
        \item Architecture monorepo avec Turbo
        \item TypeScript pour la robustesse du code
        \item Tests end-to-end avec Playwright
        \item Git workflow avec protection des secrets
        \item Documentation continue
    \end{itemize}
\end{frame}

\begin{frame}{Remerciements}
    \begin{center}
        \Huge \textcolor{primaryblue}{Merci de votre attention !}
        
        \vspace{1cm}
        
        \Large \textbf{SharingBot}
        
        \normalsize Assistant Conversationnel Intelligent avec RAG
        
        \vspace{1cm}
        
        \textbf{Questions ?}
        
        \vspace{0.5cm}
        
        \small
        \href{https://github.com/votre-repo/SharingBotV3}{GitHub: github.com/votre-repo/SharingBotV3}
        
        \href{mailto:votre.email@exemple.com}{Contact: votre.email@exemple.com}
    \end{center}
\end{frame}

\begin{frame}{Annexe : Ressources}
    \textbf{Liens Utiles}
    \begin{itemize}
        \item Next.js : \url{https://nextjs.org}
        \item Supabase : \url{https://supabase.com}
        \item OpenAI : \url{https://openai.com}
        \item Shadcn UI : \url{https://ui.shadcn.com}
        \item React Speech Recognition : \url{https://github.com/JamesBrill/react-speech-recognition}
    \end{itemize}
    
    \vspace{0.5cm}
    
    \textbf{Documentation Technique}
    \begin{itemize}
        \item README.md du projet
        \item Documentation API dans \texttt{/docs}
        \item Guide de déploiement
        \item Guide de contribution
    \end{itemize}
\end{frame}

\end{document}
